% 第 1 章：架构总览与数据流
% 本章文件由 \input 引入，不含导言区。
\section{架构总览与数据流}\label{sec:overview}

\subsection{从工程模型到求解结果的流水线}
平台潮流计算遵循统一的五段式流水线（工程语义层 $\to$ 校验 $\to$ 投影 $\to$
装配求解 $\to$ 结果回投影），全部潮流入口共用同一条主链：

\begin{figure}[H]
\centering
\begin{tikzpicture}[font=\footnotesize,
  box/.style={draw,rounded corners,minimum height=1.0cm,align=center,inner sep=4pt},
  arr/.style={->,thick}]
  \node[box] (rich) at (0,0) {\texttt{HybridPowerSystem}\\工程语义层（rich）};
  \node[box] (valid) at (4.1,0) {四级校验\\Basic/Electrical/\\SolverReady/Strict};
  \node[box] (proj) at (8.2,0) {canonical 投影\\开关收缩·零阻抗合并\\设备聚合};
  \node[box] (data) at (12.3,0) {\texttt{SolverData} 装配\\$Y_{bus}$·$G_{dc}$·注入};
  \node[box] (solve) at (4.1,-2.3) {求解/分析\\Newton·FDPF·BFS·…};
  \node[box] (back) at (8.2,-2.3) {结果回投影\\\srcpath{unproject\_pf\_result}\\支路功率回算};
  \draw[arr] (rich) -- (valid);
  \draw[arr] (valid) -- (proj);
  \draw[arr] (proj) -- (data);
  \draw[arr] (data.south) |- (solve.west);
  \draw[arr] (solve) -- (back);
\end{tikzpicture}
\caption{潮流计算主流水线（映射 \fld{bus\_merge\_map} 随 SolverData 携带，
供结果归因回原始组件编号空间）}
\end{figure}

投影与回投影的不变量（\srcpath{docs/developer/projection\_and\_results.md} 与
\srcpath{docs/archive/reference/technical\_notebook/sections/02e\_projection\_result\_contracts.tex}，
代码实现 \srcpath{src/model/network\_utils.cpp}）：强度量（电压、相角）穿越合并
\textbf{广播}；广延量（功率、切负荷、能量）必须按参与语义\textbf{显式拆分}，
绝不盲广播；合并/死岛删除映射是复合而非覆盖；恢复等级
（strong/approximate/audit-only/unsupported）是结果契约的一部分。

\subsection{生产主调用链}
统一潮流主入口 \srcpath{hacdcpf::solve\_power\_flow}
（\srcpath{src/api/hacdcpf.cpp:solve\_power\_flow}）的执行顺序如下，每一步均注明源码位置：

\begin{enumerate}
  \item \textbf{参考母线资格校验}：\srcpath{validate\_reference\_bus\_eligibility}
        （src/api/hacdcpf.cpp:solve\_power\_flow），隔离/失电母线不得作 slack；
  \item \textbf{换流器协调预校核}（可选，默认关）：
        \srcpath{powerflow::evaluate\_converter\_coordination}（:solve\_power\_flow），
        按 DC 电压岛汇总电压源与功率平衡裕度，Error/Fatal 级问题阻断求解
        （规则目录见第~\ref{sec:hybrid}~章）；
  \item \textbf{图拓扑预检}（:solve\_power\_flow）：无任何带 slack 的合法 AC 岛时直接返回失败；
  \item \textbf{SolverData 获取与缓存}：\srcpath{get\_cached\_solver\_data}（调用点 :solve\_power\_flow，
        thread\_local 缓存，指针+存储视图+哈希签名三重匹配，:get\_cached\_solver\_data）；
        未命中时 \srcpath{make\_solver\_data} 重建
        （\srcpath{src/power\_flow/solver\_data.cpp:make\_solver\_data}）；
  \item \textbf{ZIP 权重应用}：\srcpath{apply\_zip\_weights}（\srcpath{src/api/hacdcpf.cpp:apply\_zip\_weights}）；
  \item \textbf{Newton 求解}：\srcpath{engine::NewtonSolver::solve}（调用点 :solve\_power\_flow，
        实现 \srcpath{src/power\_flow/newton\_solver.cpp:solve}）；
  \item \textbf{派生结果回算}：\srcpath{populate\_derived\_results}（\srcpath{src/api/hacdcpf.cpp:populate\_derived\_results}），
        含支路潮流 \srcpath{compute\_branch\_flows} 与换流器传输功率/损耗/限值校核；
  \item \textbf{换流器物理限值强制}（可选，默认关）：
        \fld{enforce\_converter\_physical\_limits} 开启时（\srcpath{src/api/hacdcpf.cpp:solve\_power\_flow}），
        违反换流器不等式的 Newton 根被判物理不可行（详见第~\ref{sec:hybrid}~章）；
  \item \textbf{结果回投影}：\srcpath{unproject\_pf\_result}（\srcpath{src/api/hacdcpf.cpp:unproject\_pf\_result}），
        把 canonical 解广播/归因回原始母线与支路编号空间，并恢复被合并开关/
        断路器的端子潮流。
\end{enumerate}

\subsection{公共 API 入口族}
全部潮流相关公共入口（\srcpath{src/api/hacdcpf.cpp} 与
\srcpath{include/hacdcpf/api/hacdcpf.hpp}）汇总如下；GUI 后端
（\srcpath{tests/run\_gui\_server.cpp}）的 \texttt{method} 分发与之对应
（\texttt{ac\_newton}/\allowbreak\texttt{pure\_ac}/\allowbreak\texttt{dc}/\allowbreak\texttt{fdpf}/
\texttt{adaptive}/\allowbreak\texttt{islanded}/\allowbreak\texttt{distributed\_slack}/
\texttt{three\_phase}/\texttt{three\_phase\_hybrid} 等）：

\begin{center}
\footnotesize
\begin{tabular}{@{}L{6.2cm}L{3.2cm}L{5.6cm}@{}}
\toprule
入口 & 求解器/路径 & 说明\\
\midrule
\srcpath{solve\_power\_flow} & 统一 Newton & 混合 AC/DC 主入口\\
\srcpath{powerflow::solve\_ac} & $\to$ \srcpath{solve\_power\_flow} & 薄包装（\srcpath{src/power\_flow/ac.cpp:solve\_ac}）\\
\srcpath{powerflow::solve\_hybrid} & $\to$ \srcpath{solve\_power\_flow} & 薄包装（\srcpath{src/power\_flow/hybrid.cpp}）\\
\srcpath{solve\_power\_flow\_fdpf} & \srcpath{FDPFSolver} & 快速解耦，纯 AC\\
\srcpath{solve\_ac\_linearized\_dc} & $B\theta=P$ & 线性化角度潮流（\srcpath{src/power\_flow/ac\_linearized\_pf.cpp:solve\_ac\_linearized\_dc}）\\
\srcpath{solve\_dc\_power\_flow} & \srcpath{DCSolver} & DC 网络电压 Newton（\srcpath{src/power\_flow/dc\_solver.cpp:solve}）\\
\srcpath{solve\_power\_flow\_adaptive} & \srcpath{AdaptiveSolver} & 分岛并行（\srcpath{src/power\_flow/adaptive\_solver.cpp:solve}）\\
\srcpath{solve\_power\_flow\_distributed\_slack}(\srcpath{\_full}) & \srcpath{DistributedSlackSolver} & 分布式松弛两变体\\
SolverHandle 族 & \srcpath{make\_solver\_handle} 等 & 符号分解/模式跨多次求解复用（时序仿真用）\\
\srcpath{analysis::solve\_distribution\_pf} & BFS 前推回代 & 辐射状配网（\srcpath{src/power\_flow/distribution\_power\_flow.cpp:solve\_distribution\_pf}）\\
\srcpath{analysis::solve\_three\_phase\_nr} & 相域 NR & 三相不平衡（\srcpath{src/power\_flow/three\_phase\_nr.cpp:solve\_three\_phase\_nr\_snapshot}）\\
\srcpath{solve\_three\_phase\_hybrid\_pf} & 相域 AC+DC 联立 & 矩阵级 API（\srcpath{src/power\_flow/three\_phase\_hybrid.cpp:solve\_three\_phase\_hybrid\_pf}）\\
\srcpath{CpfSolver::solve} & 连续潮流 & 电压稳定（\srcpath{src/power\_flow/voltage\_stability.cpp:solve}）\\
\bottomrule
\end{tabular}
\end{center}

\subsection{求解器抽象层与工厂}
抽象接口 \srcpath{IPowerFlowSolver}（\srcpath{include/hacdcpf/power\_flow/solver\_interface.hpp:IPowerFlowSolver}）
声明 \texttt{solve(problem)}/\texttt{name()} 两方法及 6 个子类
（\texttt{NewtonHybridSolver}/\allowbreak\texttt{FDPFSolver}/\allowbreak\texttt{DCPowerFlowSolver} 等，:NewtonHybridSolver、:FDPFSolver、:DCPowerFlowSolver、:AdaptiveHybridSolver、:ContinuationPowerFlowSolver、:NewtonKrylovSolver）；
问题描述符 \srcpath{PowerFlowProblem}（\srcpath{include/hacdcpf/power\_flow/power\_flow\_problem.hpp:PowerFlowProblem}）携带
SolverData、选项与规模摘要，工厂 \srcpath{build\_power\_flow\_problem}
（声明 \srcpath{include/hacdcpf/power\_flow/power\_flow\_problem.hpp:build\_power\_flow\_problem}，实现
\srcpath{src/power\_flow/solver\_interface.cpp:build\_power\_flow\_problem}）装配。
\srcpath{solver\_factory.hpp} 提供 \texttt{PowerFlowMethod} 枚举（含
\texttt{Continuation}）与显式 \texttt{create(PowerFlowMethod)}；六个子类的 \texttt{solve}
与工厂实现同集中于新增的 \srcpath{src/power\_flow/solver\_interface.cpp}
（:solve）。算法选择须经 \texttt{PowerFlowMethod}\textbf{显式}进行；
\texttt{PowerFlowOptions} 只配置已选方法，不再提供会被忽略的工厂重载。
如实说明：抽象层现已完整接线（\texttt{NewtonHybridSolver::\allowbreak solve}
委托 \srcpath{powerflow::NewtonSolver}，\srcpath{src/power\_flow/solver\_interface.cpp:solve}），
但生产门面 \srcpath{solve\_power\_flow} 仍不经过该抽象层，直接持有
thread\_local 的 \srcpath{engine::NewtonSolver}（src/api/hacdcpf.cpp:solve\_power\_flow；
规范头 \srcpath{include/hacdcpf/engine/solver/native/nle/newton\_solver.hpp}；
\srcpath{power\_flow/newton\_solver.hpp} 与 \srcpath{power\_flow/solvers/newton\_solver.hpp}
均为转发头）。

\subsection{零阻抗合并对潮流装配的影响}
母线合并发生在投影层（装配之前）：
\srcpath{merge\_zero\_impedance\_buses\_impl}（\srcpath{src/model/network\_utils.cpp:merge\_zero\_impedance\_buses\_impl}）
以阈值 \fld{impedance\_threshold}=$10^{-4}$（\srcpath{include/hacdcpf/projection/project\_to\_canonical.hpp:ProjectionOptions}）
合并近零阻抗支路两端母线；闭合断路器先折叠为近零阻抗支路再被合并
（src/model/network\_utils.cpp:add\_equivalent\_branch\_from\_switch）。因此 SolverData 中只剩 canonical 母线，
$Y_{bus}$ 装配时 $|z|=0$ 支路直接跳过（\srcpath{src/power\_flow/admittance\_builder.cpp:build\_admittance\_matrix}）。
合并映射 \fld{bus\_merge\_map} 随 SolverData 携带，求解后回投影恢复原始编号空间
的结果与开关端子潮流。

\subsection{缓存体系}
两层缓存避免重复装配与重复符号分解：
\begin{itemize}
  \item \textbf{API 层 SolverDataCache}（thread\_local，src/api/hacdcpf.cpp:SolverDataCache）：
        同一 \texttt{HybridPowerSystem} 对象且 \fld{loss\_model} 不变时直接复用；
  \item \textbf{求解器层 PatternCache}（规范头
        \srcpath{include/hacdcpf/engine/solver/native/nle/newton\_solver.hpp}；
        \srcpath{include/hacdcpf/power\_flow/newton\_solver.hpp} 为转发头）：
        缓存 \texttt{JacobianContext}、\texttt{JacobianPattern} 与稀疏线性求解器实例，
        键含数据指针、\fld{build\_id}（原子递增，缓存失效标记）与 Ybus/Gdc 存储地址；
        稀疏模式只在缓存失效时 \texttt{analyzePattern} 一次，之后每次迭代仅数值分解；
  \item \textbf{求解器状态工作区 SolverWorkspace}
        （\srcpath{include/hacdcpf/power\_flow/workspace.hpp:SolverWorkspace}，
        实现 \srcpath{src/power\_flow/workspace.cpp}）：已不再是空脚手架——
        \texttt{prepare\_state}/\texttt{prepare\_equations}/\texttt{reset\_iteration}
        在维度不变时保留 Eigen 分配、仅重置数值。facade 的 thread\_local
        \texttt{NewtonSolver} 跨普通求解复用同一工作区
        （成员 \fld{workspace\_}/\fld{workspace\_in\_use\_}，engine 头 :NewtonSolver；
        src/power\_flow/newton\_solver.cpp:solve 的占用守卫使同伦递归的内层求解改用局部工作区，
        避免覆盖外层状态），状态与方程向量经 \texttt{prepare\_state}（:solve）/
        \texttt{prepare\_equations}（:solve）挂接。
\end{itemize}
时序生产模拟（UC$\to$OPF$\to$PF 流水线）经 SolverHandle 族在整个时间轴上复用
同一求解器实例，拓扑仅在 \fld{in\_service} 变化时重建
（\srcpath{docs/theory/time\_series\_power\_flow\_models.md}）。
